Per la paret que teniu al darrere de l'aula on feu l'examen, entren protons amb una trajectòria horitzontal i a una velocitat $\vec{v}_\mathrm{p+} = 2,00\times 10^{6}\,\vec{\imath}\ \mathrm{m/s}$. Dins l'aula hi ha un camp magnètic també horitzontal el valor del qual és $\vec{B} = 0,500\,\vec{\jmath}\ \mathrm{T}$. Determineu:

\begin{apartat}{1}
La força causada pel camp magnètic que actua sobre els protons quan entren en la zona on hi ha aquest camp magnètic.
\end{apartat}

\begin{apartat}{1}
El radi de la trajectòria circular dels protons dins l'aula i indiqueu si aquests protons impactaran contra les persones que estan assegudes a l'aula.
\end{apartat}

\dades{%
  Càrrega del protó: $1,60\times 10^{-19}\ \mathrm{C}$\\
  Massa del protó: $1,67\times 10^{-27}\ \mathrm{kg}$}
\nota{Negligiu el pes del protó.}
